\documentclass[11pt]{article}
\usepackage[margin=2.6cm]{geometry}
\usepackage{mathpazo}
\usepackage{amssymb}
\usepackage{microtype}
\usepackage{booktabs}
\usepackage{graphicx}
\usepackage{float}
\usepackage{xcolor}
\usepackage{titlesec}
\usepackage{caption}
\usepackage{listings}
\usepackage[hidelinks]{hyperref}
\emergencystretch=3em
\graphicspath{{./}{../report_full250/}}

\definecolor{accent}{HTML}{1F4E79}
\definecolor{good}{HTML}{1A7F37}
\definecolor{bad}{HTML}{B42318}
\definecolor{qwen}{HTML}{1F4E79}
\definecolor{gem}{HTML}{C9720B}
\titleformat{\section}{\Large\bfseries\color{accent}}{\thesection}{0.8em}{}
\titleformat{\subsection}{\large\bfseries\color{accent!80!black}}{\thesubsection}{0.8em}{}
\captionsetup{font=small,labelfont={bf,color=accent}}
\newcommand{\armA}{\textcolor{bad}{arm A}}
\newcommand{\armB}{\textcolor{accent}{arm B}}
\newcommand{\partrule}[2]{%
  \section*{#1}%
  \noindent{\large\color{accent!70!black}#2}\par\bigskip}

\lstdefinestyle{json}{
  basicstyle=\ttfamily\footnotesize,
  backgroundcolor=\color{gray!8},
  frame=single, rulecolor=\color{gray!40}, framerule=0.4pt,
  breaklines=true, columns=fullflexible, keepspaces=true,
  showstringspaces=false, aboveskip=8pt, belowskip=8pt}
\lstset{style=json}

\newsavebox{\eliBox}
\newenvironment{eli5}
  {\par\medskip\noindent\begin{lrbox}{\eliBox}\begin{minipage}{\dimexpr\linewidth-2\fboxsep-2\fboxrule}\small}
  {\end{minipage}\end{lrbox}\noindent\fcolorbox{accent!40}{accent!4}{\usebox{\eliBox}}\par\medskip}

%% Held-out qwen test numbers — filled from test_metrics_qwen.json when the eval lands.
\newcommand{\qTseedS}{0.767} % seed test soft
\newcommand{\qTseedD}{48.3} % seed test discrete %
\newcommand{\qTaS}{0.833}    % arm A winner test soft
\newcommand{\qTaD}{59.0}    % arm A winner test discrete %
\newcommand{\qTbS}{0.816}    % arm B winner test soft
\newcommand{\qTbD}{56.7}    % arm B winner test discrete %
\newcommand{\qTgS}{0.827}    % arm B gen2 test soft
\newcommand{\qTgD}{61.0}    % arm B gen2 test discrete %

\title{\color{accent}\bfseries Tool-Aware DAG-Diff Mutation for Retrieval Chains\\
\large Two paired A/Bs on HoVer multi-hop retrieval: does the structured diff still\\
earn its keep as the mutator scales from a mid-tier to a top-tier reasoning model?\\
\large (single-parent; mutators \texttt{gemini-3.5-flash} and \texttt{Qwen3-235B-A22B-Thinking}; full 250-mutant runs)}
\author{GigaEvo}
\date{2026-07-05}

\begin{document}
\maketitle

\begin{abstract}
\noindent
Two paired A/B runs on HoVer multi-hop retrieval compare \emph{free full-document rewriting}
(\armA) against a \emph{schema-compiled, tool-aware slot diff} (\armB) as the mutation operator.
The executor is held fixed; only the mutator scales: \texttt{gemini-3.5-flash} (Part~I) and
\texttt{Qwen3-235B-A22B-Thinking} (Part~II). A first attempt produced a spectacular --- and
spurious --- headline: ``38\% of a strong mutator's rewrites are malformed JSON.'' An audit traced
almost all of it to a harness bug: the structured-output schema demanded Python while the prompt
demanded JSON. We fixed the schema and re-ran both \armA{} runs, keeping the originals as a study
in structured-output hygiene (\S\ref{sec:artifact}). The honest post-fix picture: free rewriting
under a strong mutator is \emph{not} a validity disaster (96.7\% / 85.8\% valid), and peak fitness
is statistically tied in both pairs --- the same seed re-scored four times spans 0.732--0.753, a
noise band wider than either best-fitness gap. Three advantages survive at both mutator strengths.
\textbf{Validity:} the diff persisted \emph{zero} malformed children in both of its runs; every
residual free-rewrite child that persisted malformed is a \emph{truncated long emission}
(tail-clipped on gemini, 3.3\%; head-clipped on the local Qwen serving stack, 13.8\%) --- a
fragility of whole-document re-emission that the diff's short, slot-local payloads avoided on
both stacks. \textbf{Search efficiency:} in these runs the diff's never-beaten winner appeared at
31\% / 24\% of the budget versus 52\% / 94\%, so retrospective cost-to-winner falls \textbf{36\%}
(Part~I) and \textbf{70\%} (Part~II) at parity cost per valid program --- point estimates from one
replicate per arm. \textbf{Legibility:} both diff-arm winners sit on short, typed lineages whose
edits --- including a final prompt-only, single-field mutation --- can be read verbatim; the
rewrite arm has no operator-native diff to inspect. Verdict: default to the diff for wire-JSON
genomes at every mutator strength tested; after harness hygiene its dividend is not ``strong
models can't write JSON'' but \emph{earlier, cheaper, auditable winners with zero persisted
malformed output}.
\end{abstract}

\section{Motivation}
On the CARL chain tasks the evolved ``program'' is not Python but a wire-JSON document
describing a DAG of retrieval and LLM steps (\texttt{ReasoningChain.to\_dict} plus platform
extras). The stock mutation path (\texttt{mutation=llm\_rewrite}) asks the LLM to re-emit the
\emph{whole} document as free text: every emission re-risks the entire syntax, and a malformed
child is only discovered \emph{after} it has been persisted and (if the parse survives
extraction) has consumed a full evaluation. Two prior A/Bs established the structured-diff
alternative on the \emph{summarizer} task, whose chains are LLM-only. This report answers three
questions those runs left open; the third is the reason for the second, cross-model run:
\begin{itemize}
  \item \textbf{Does the diff generalise to tool-bearing genomes?} HoVer chains interleave
  \texttt{tool} steps (\texttt{retrieve}, \texttt{retrieve\_deep}) with LLM steps. The
  original \texttt{ChainDagDiff} was LLM-step-only and hard-rejected tool steps; \armB{}
  exercises a new vocabulary, \texttt{AllowedToolChainChanges}, that admits tool steps as
  first-class slots.
  \item \textbf{Does the guarantee still matter when the mutator is \emph{strong}?} The prior
  run used a weak non-reasoning 8B model, where free-form JSON was expected to break. Part~I's
  mutator is \texttt{gemini-3.5-flash} with \texttt{reasoning.effort:high} --- a capable model
  that should rarely emit malformed JSON.
  \item \textbf{Does the advantage \emph{survive} an even stronger mutator?} Part~II re-runs the
  identical A/B with \texttt{Qwen3-235B-A22B-Thinking} --- a frontier-scale reasoning model. If
  the diff's edge were merely ``the mutator is not good enough at JSON,'' a stronger mutator
  should erase it. The cross-model comparison (\S\ref{sec:synthesis}) separates the advantages
  that \emph{shrink} with mutator strength from the ones that \emph{hold}.
\end{itemize}
A fourth finding was unplanned. The first \armA{} runs under both mutators were invalidated by a
schema/prompt contradiction in the mutation harness; \S\ref{sec:artifact} documents the bug, the
decomposition of the spurious ``malformed'' children, and what the episode teaches about
structured-output hygiene. All \armA{} numbers in Parts I--II come from the fixed re-runs.

\section{Setup}
\label{sec:setup}
Two \emph{paired} A/Bs, each two arms run from scratch, arms identical except the mutation
operator; the two pairs identical except the mutator model (\texttt{launch\_arm.sh},
\texttt{launch\_arm\_qwen.sh}):
\begin{itemize}
  \item \textbf{Task:} \texttt{problems/chains/hover/full7}, \texttt{pipeline=program\_is\_json}.
  Full-chain mode: no frozen steps, no topology matching; the child may freely change step
  count/types/dependencies within \texttt{FULL\_CHAIN\_CONFIG} (\texttt{max\_steps}=7). Fitness
  is soft retrieval coverage --- for each of 300 HoVer claims, the fraction of gold
  supporting-fact article titles found across \emph{all} tool-step outputs (higher better).
  \item \textbf{Mutators (the only cross-pair difference):}
  \textbf{Part~I} --- \texttt{google/gemini-3.5-flash} via OpenRouter, with
  \texttt{reasoning.effort:high} and \texttt{service\_tier:flex}.
  \textbf{Part~II} --- \texttt{Qwen3-235B-A22B-Thinking-2507} served by the local LiteLLM
  proxy. Within each pair the mutator config is \emph{identical} across arms.
  \item \textbf{Executor:} \texttt{Qwen/Qwen3-8B} via the LiteLLM proxy, run step-batched over the
  300-sample set --- \emph{the same executor in both pairs}. Holding it fixed removes the executor
  as a designed source of variation; stochastic scoring noise and serving-stack state remain
  (see the caveats).
  \item \textbf{Config:} \texttt{memory=none}, \texttt{num\_parents=1},
  \texttt{max\_mutants=250}, one 7-step retrieval seed, disk storage. \armA{} runs the stock
  full wire-JSON rewrite
  (\texttt{mutation=llm\_rewrite}). \armB{} runs the \texttt{StructuredDiffMutationOperator}
  driving the \texttt{AllowedToolChainChanges} grammar
  (\texttt{mutation=carl\_with\_retrieval\_tools}).
\end{itemize}

\noindent\textbf{Single parent.} With \texttt{num\_parents=1} the diff grammar collapses to its
single-parent form (\texttt{base\_parent} is one letter; every \texttt{keep}-id enum ranges
over one parent's steps). Crossover is unavailable this run --- a deliberate simplification,
two-parent tool-crossover left as prior work.

\medskip\noindent\textbf{Provenance and timing.} The \armA{} runs reported in Parts I--II are
re-runs, launched after the harness fix of \S\ref{sec:artifact}; the \armB{} runs predate them by
several days. Configs, seed, and executor are unchanged in between, but the arms of each pair did
not run concurrently, so shared proxy and cluster load over those days are uncontrolled --- a
caveat we accept and disclose rather than re-run \armB.

\medskip\noindent\textbf{Fitness is a noisy draw.} The executor samples at temperature 0.6
(top-p 0.95), so a program's fitness is stochastic. The same seed chain, re-scored once per run,
came back 0.732 / 0.753 (Part~I A/B) and 0.734 / 0.743 (Part~II A/B) --- $0.741 \pm 0.010$
across the four draws, a \textbf{0.021 spread with identical genomes}. Best-fitness gaps at or below that band
(both pairs, as it turns out) should be read as ties. The band bounds noise on
\emph{single-program} statistics --- a best child, a seed re-score; arm \emph{means} over
$\sim$200+ children carry standard errors an order of magnitude smaller and are read against the
per-arm spread instead (the mean~$\pm$~SD rows in the fitness tables). Best-of-$N$ selection bias matters too:
an arm with more valid children takes more draws at the maximum. Part~I is nearly balanced
(234 vs.\ 239 valid); in Part~II the \emph{rewrite} arm has fewer draws (211 vs.\ 245) yet still
posts the higher best --- so draw-count bias does not explain its edge, though the noise band does.

\section{The tool-aware diff language}
\label{sec:language}
The child chain is a shelf of numbered slots; the LLM fills them left-to-right. What is new
here versus the summarizer diff is that a slot may hold a \texttt{tool} step, not only an LLM
step: a kept tool step carries its tool name and arguments, and the grammar admits editing them.

\begin{eli5}
Picture the child chain as a row of numbered boxes. Into each box you put either a
\emph{photocopy} of a step from the parent --- named by its label (\texttt{a1}, \texttt{a2},
\dots), optionally with small corrections --- or a \emph{brand-new} step. Some parent steps are
``go-look-it-up'' \emph{tool} steps (retrieve passages); those are photocopied and corrected the
same way as thinking steps. Each box lists which \emph{earlier} boxes it reads from, and the
form for box $k$ only has checkboxes for boxes $1..k{-}1$ --- so pointing forward or at yourself
is not against the rules, there is simply nowhere on the form to write it. Leftover boxes stay
empty (no gaps), and the last filled box produces the graded answer.
\end{eli5}
\noindent
The form is a JSON schema compiled fresh per mutation from the actual parent and shipped to the
provider as the \texttt{json\_schema} structured-output grammar, so a violating token can never
be sampled. A pydantic model re-checks the payload as a backstop, and the applier round-trips
the result through \texttt{ReasoningChain} as a final assertion. Everything an unconstrained rewriter
routinely gets wrong --- self/forward dependencies, cycles, keeping a nonexistent step,
over-long chains, invented fields, \emph{malformed JSON} --- is unrepresentable by construction.
Only step-contiguity is left to a validator, with one disclosed subtlety: a payload whose filled
slots are non-contiguous is not rejected but \emph{compacted} (filled slots slide left, order
preserved). This silent-repair semantic went unlogged in the runs below; the current code emits
a warning whenever it fires, so future runs can bound its frequency. A rejection costs one LLM
call: nothing persisted, nothing evaluated.

\partrule{Part I --- mid-tier mutator: \texttt{gemini-3.5-flash}}{full 250-mutant paired run, \texttt{reasoning.effort:high}}

\section{Results (gemini mutator)}
\label{sec:results}

\subsection{Validity funnel --- a residual truncation tax, not a catastrophe}
\begin{table}[H]\centering
\begin{tabular}{lcc}
\toprule
full 250-mutant run & \armA{} (rewrite) & \armB{} (tool diff) \\
\midrule
mutation LLM calls (incl.\ retries) & 276 & 259 \\
pre-persist generation failures & 24 (no code extracted) & 4 (2 schema, 2 call) \\
children persisted & 252 & 254 \\
\quad still executing at stop (excluded) & 10 & 12 \\
\midrule
\textbf{completed evaluations} & \textbf{242} & \textbf{242} \\
\quad valid & 234 & \textbf{239} \\
\quad invalid & 8 & 3 \\
\qquad of which malformed wire-JSON & \textcolor{bad}{8 (3.3\%)} & \textbf{\textcolor{good}{0}} \\
\qquad other (non-genome) & 0 & 3$^\ast$ \\
\textbf{valid yield among completed} & 96.7\% & \textbf{\textcolor{good}{98.8\%}} \\
\bottomrule
\end{tabular}
\caption{With the harness fixed, the free-form path is no longer a validity disaster: 8 of 242
completed children (3.3\%) are malformed wire-JSON, versus the pre-fix run's 38\%
(\S\ref{sec:artifact}). All 8 are \emph{tail-truncated}: each document opens correctly and dies
mid-object at its final character --- consistent with the emission being cut off (e.g.\ an output
ceiling binding after a long reasoning phase), not with syntax confusion. \armB{} persisted zero
malformed children; $^\ast$its 3 completed-invalids are evaluator-infrastructure crashes
(\texttt{exit=1}) whose genomes parse fine, so \armB's genome-defect count is 0. \armA's 24
pre-persist failures (8.7\% of calls) returned no extractable document at all --- a wasted call
each, but nothing persisted or evaluated. (Funnel rows need not close exactly --- e.g.\
$259-4=255$ calls vs 254 persisted in \armB{}: the odd residual call produced output that was
never persisted and is unclassifiable from the logs.)}
\end{table}
\begin{figure}[H]\centering
\includegraphics[width=\textwidth]{ab_overview.png}
\caption{Left: per-child soft-coverage fitness (dots) and best-so-far (line); invalid children
are drawn as crosses on the floor line --- now sparse in \emph{both} arms. \armB's best-so-far
reaches 0.858 vs \armA's 0.843. Right: the validity funnel over the full run --- both persist
$\sim$250, \armA{} keeps 234 valid to \armB's 239.}
\end{figure}

\noindent\textbf{What still breaks, and why it matters.} The pre-fix story (``a strong model
emits malformed JSON in 38\% of its whole-document rewrites'') is dead --- that was our harness
(\S\ref{sec:artifact}). What survives is smaller but structural: a whole-document rewrite is a
single $\sim$2--5k-token emission, and \emph{some} of those get truncated in flight. The diff
never ships a payload that long --- kept slots are named, not re-serialised --- so a truncation
of that depth has nothing to bite on; and were the outer response cut anyway, the
grammar-checked parse fails \emph{pre-persist} rather than persisting a corrupt genome.
Part~II tests whether this residue also appears on a different serving stack (it does,
spectacularly, from the other end of the document).

\subsection{Fitness --- \armB{} leads, inside the noise band}
\begin{table}[H]\centering
\begin{tabular}{lcc}
\toprule
completed children & \armA{} (rewrite) & \armB{} (tool diff) \\
\midrule
best child & 0.843 & \textbf{\textcolor{good}{0.858}} \\
child mean $\pm$ SD & $0.782 \pm 0.031$ & \textbf{\textcolor{good}{$0.800 \pm 0.045$}} \\
child median & 0.789 & \textbf{\textcolor{good}{0.802}} \\
seed (executor re-score) & 0.732 & 0.753 \\
\bottomrule
\end{tabular}
\caption{\armB{} leads all central measures. The best-child margin (0.015) sits inside the 0.021
seed re-score spread --- the same genome scored 0.732 in one run and 0.753 in the other --- and
reads as a tie. The mean gap ($+0.018$; per-arm standard errors $\approx$0.002--0.003) is larger
than per-evaluation sampling noise alone, but with one replicate per arm and non-concurrent
execution we read it as a consistent \armB-ward lean, not an established win.}
\end{table}

\subsection{Held-out test coverage vs.\ the CARL feedback-condition baselines}
\label{sec:benchcmp}
The four inference-time feedback conditions (\emph{without-feedback}, \emph{self-reason},
\emph{metric-aware}, \emph{data-aware}) are an \textbf{executor} axis, orthogonal to this
report's \textbf{mutator} axis. Our runs hold the executor fixed at \texttt{feedback\_mode=none}
(\emph{without-feedback}: single-pass, no retry) and vary only the mutation operator. To place
our winners on the same HoVer test set, Table~\ref{tab:benchcmp} reproduces the feedback-condition
sweep (mutator \texttt{Qwen3-235B-A22B-Thinking}) and adds \emph{both} mutator pairs' best
programs. Metric is \textbf{discrete coverage} ($1$ iff every gold title
retrieved), $n=300$ test claims, base $=$ seed chain $\to$ best $=$ evolved winner.

\begin{table}[H]\centering\small
\begin{tabular}{lllcl}
\toprule
method & mutator & feedback & test disc.\ base$\to$best & source \\
\midrule
prompt & Qwen3-235B-Think & ---          & $43.3 \to 49.0$ & prior table$^\dagger$ \\
GEPA   & Qwen3-235B-Think & ---          & $43.3 \to 46.0$ & prior table$^\dagger$ \\
CARL   & Qwen3-235B-Think & none         & $48.7 \to 56.0$ & prior table$^\dagger$ \\
CARL   & Qwen3-235B-Think & self-reason  & $46.3 \to 57.0$ & prior table$^\dagger$ \\
CARL   & Qwen3-235B-Think & metric-aware & $46.3 \to 60.7$ & prior table$^\dagger$ \\
CARL   & Qwen3-235B-Think & data-aware   & $68.3 \to 68.7$ & prior table$^\dagger$ \\
\midrule
\armA{} (free rewrite)   & gemini-3.5-flash & none & $44.0 \pm 2.9 \to \textbf{67.3} \pm 2.7$ & \textbf{Part I} \\
\armB{} (CARL tool diff) & gemini-3.5-flash & none & $44.0 \pm 2.9 \to \textbf{\textcolor{good}{68.0}} \pm 2.7$ & \textbf{Part I} \\
\addlinespace
\armA{} (free rewrite)   & Qwen3-235B-Think & none & $\qTseedD \pm 2.9 \to \qTaD \pm 2.8$ & \textbf{Part II} \\
\armB{} (CARL tool diff) & Qwen3-235B-Think & none & $\qTseedD \pm 2.9 \to \qTbD \pm 2.9$ & \textbf{Part II} \\
\bottomrule
\end{tabular}
\caption{Held-out HoVer discrete coverage ($n=300$). Our rows carry the binomial
claim-sampling $\pm$SE on each value, with executor re-draw noise on top; the prior-table rows
are user-supplied point values with no dispersion available. \textbf{Top block:} the feedback-condition
sweep under the \texttt{Qwen3-235B-Thinking} mutator (user-supplied; \emph{prompt} and
\emph{GEPA} are the sweep's prompt-optimisation baselines). \textbf{Middle block:}
Part~I's two winners under the \texttt{gemini-3.5-flash} mutator, both at \texttt{feedback=none}
--- a near-tie (68.0 vs 67.3, within the $\sim$3.7-pt single-draw noise of
Table~\ref{tab:lineagetest}); both sit $\sim$11 points above the prior sweep's without-feedback
best (56.0), though that comparison crosses evaluations ($^\dagger$). \textbf{Bottom block:}
Part~II's winners under the \emph{same} mutator as the prior sweep --- the closest available
comparison, with modest margins (59.0 / 56.7 vs 56.0, inside the noise floor;
\S\ref{sec:benchcmpq}). Our seed re-scores inside the prior table's base cluster
($43.3$--$48.7$), consistent with --- though not proof of --- a shared metric and split.
$^\dagger$Prior rows are user-supplied; their column semantics (Val vs Test) and the identity of
their ``CARL'' operator are our best-judgment reading, so treat cross-block comparisons as
contextual, not conclusive.}
\label{tab:benchcmp}
\end{table}

\subsection{Token cost --- full reasoning split, near-parity totals}
\label{sec:tokens}
The pre-fix run could not compare reasoning tokens: \armA's emitter silently dropped the field
\texttt{tokens\_reasoning}. That instrumentation bug is fixed, and the re-run gives the
first complete cross-arm split under a reasoning mutator.
\begin{table}[H]\centering\small
\begin{tabular}{lccc}
\toprule
mutation side, full run (gemini) & \armA{} & \armB{} & \armB{} vs \armA{} \\
\midrule
tokens in & 1{,}821{,}207 & 1{,}934{,}479 & $+6.2\%$ \\
tokens out (total) & 1{,}890{,}914 & 1{,}882{,}208 & $-0.5\%$ \\
\quad of which reasoning & 1{,}502{,}356 (79\%) & 1{,}639{,}000 (87\%) & $+9.1\%$ \\
\quad \textbf{content (out $-$ reasoning)} & \textbf{388{,}558} & \textbf{243{,}208} & \textbf{$-37\%$} \\
LLM calls (incl.\ retries) & 276 & 259 & $-6\%$ \\
\bottomrule
\end{tabular}
\caption{Total output is a wash: gemini reasons hard on both arms (79\% and 87\% of output), and
\armB's heavier thinking almost exactly offsets its lighter payload. The \emph{content} row is
the diff's designed economy finally measured cleanly: emitting slot deltas instead of the whole
re-serialised chain costs \textbf{37\% fewer content tokens}. \armB{} pays $+6\%$ input for the
schema-bearing prompt. Net: on a mid-tier reasoning mutator, the diff's token bill is the same
--- its economy shows up in \emph{what} is generated, not how much.}
\end{table}
\begin{figure}[H]\centering
\includegraphics[width=\textwidth]{ab_tokens.png}
\caption{Left: mutation output tokens per call (dots) with a 10-call rolling mean --- the arms
overlap; \armB's small content payload plus its larger reasoning channel matches \armA's
content-heavy output. Right: cumulative tokens --- input and output curves track closely
throughout; no separation like the non-reasoning summarizer run.}
\end{figure}

\subsection{Cost/quality frontier --- the dollar that matters is cost-to-winner}
\label{sec:pareto}
Converting tokens to dollars at the cheapest OpenRouter \texttt{gemini-3.5-flash} price
(\$1.50/M in, \$9.00/M out, output inclusive of reasoning; the shared Qwen3-8B executor is
excluded) reframes the comparison. Cost is attributed by \texttt{created\_at}: each child, in
wall-clock creation order, consumes the next accepted mutation call, so ``cost to a program'' is
the cumulative mutation spend up to and including its creation. (We deliberately do \emph{not}
use the persisted \texttt{atomic\_counter} field --- it is non-monotone against creation time
here; \texttt{created\_at} and the evolution-loop \texttt{iteration} clock agree to within a
rank or two throughout.)

\begin{table}[H]\centering
\begin{tabular}{lccc}
\toprule
mutation-model \$ (OpenRouter cheapest) & \armA{} & \armB{} & \armB{} vs \armA{} \\
\midrule
total run spend & \$19.75 & \$19.84 & $+0.5\%$ \\
\quad of which reasoning channel & \$13.52 (68\%) & \$14.75 (74\%) & --- \\
\midrule
\textbf{cost to the winning program} & \$10.17 & \textbf{\textcolor{good}{\$6.46}} & \textbf{$-36\%$} \\
\quad winner reached at (chrono.\ rank) & 131/252 (52\%) & \textbf{\textcolor{good}{80/254 (31\%)}} & earlier \\
\quad winner fitness (train soft) & 0.843 & 0.858 & tie (noise band) \\
\midrule
cost per \emph{valid} program & \$0.084 & \$0.083 & parity \\
budget wasted on invalid children & \$0.65 (3.3\%) & \$0.25 (1.2\%) & --- \\
\bottomrule
\end{tabular}
\caption{Post-fix, total spend and cost-per-valid-program are parity --- the pre-fix run's
``38\% of \armA's budget wasted'' collapses to \$0.65 once the harness stops sabotaging the
rewrites. What remains is \emph{timing}: \armB's never-beaten winner exists after 31\% of the
run (\$6.46), \armA's only after 52\% (\$10.17) --- retrospectively $-36\%$ cost-to-winner.
Banking that saving online would need a stopping rule; the peak itself is a statistical tie.}
\end{table}
\begin{figure}[H]\centering
\includegraphics[width=\textwidth]{ab_pareto.png}
\caption{\textbf{Left:} cost-to-winner vs.\ best train fitness --- each arm's best-so-far
fitness as cumulative mutation dollars accrue. \armB{} sits up-and-to-the-left and owns the
frontier. \textbf{Right:} cost-to-program vs.\ held-out test discrete coverage, every evaluated
lineage program placed at its true cost-up-to-creation. \armB's connected climb reaches $68\%$
coverage by \$6.46, with the decisive \texttt{gen1}$\to$\texttt{gen2} rewrite already at
$64.3\%$ for a cumulative \textbf{\$0.69}; \armA's winner reaches $67.3\%$ at \$10.17. x-axis is
mutation-model tokens only; executor excluded.}
\end{figure}
\begin{figure}[H]\centering
\includegraphics[width=\textwidth]{ab_waste.png}
\caption{\textbf{Left:} each arm's mutation budget split (by its valid/invalid child rate) into
productive spend (green) and spend on invalid children (red) --- \$0.65 vs \$0.25, both small
post-fix. \textbf{Right:} cost per valid program, \$0.084 vs \$0.083 --- parity. The waste
asymmetry that dominated the pre-fix run now lives almost entirely in \emph{when} the winner
arrives (previous figure), not in burned evaluations.}
\end{figure}

\subsection{Structural exploration --- both preserve topology, \armB{} tighter}
\label{sec:struct}
Both arms mostly preserve the seed's 7-step structure. \armA's valid children are
$\{7\!:\!181,\ 6\!:\!32,\ 5\!:\!21\}$ (77\% at 7 steps); \armB's are
$\{7\!:\!200,\ 6\!:\!20,\ 5\!:\!18,\ 3\!:\!1\}$ (84\% at 7 steps), with a single deliberate
$7\!\to\!3$ prune. No wholesale chain-collapse on this task --- retrieval coverage rewards
keeping tool steps, so even the free rewriter keeps them --- but the diff stays marginally
closer to parent topology, consistent with edits that touch slots rather than re-imagine the
chain. (The strong-mutator pair inverts this; see \S\ref{sec:structq}.)
\begin{figure}[H]\centering
\includegraphics[width=0.72\textwidth]{ab_nsteps.png}
\caption{Chain-length distribution of valid children. Both arms concentrate at 7 steps; \armA{}
spreads a little more toward 5--6, \armB{} stays tighter at 7 with one 3-step prune.}
\end{figure}

\section{The winning chains --- convergent architecture}
\label{sec:chains}
The two arms converge on essentially the same best retrieval architecture. Both winners abandon
the seed's ``retrieve $\to$ summarise $\to$ query'' rhythm for a \textbf{deepen-and-accumulate}
chain: four \texttt{retrieve\_deep} ($k{=}10$) tool steps alternating with three
query-generation LLM steps, the summarise steps --- an information bottleneck that discarded raw
passages before the next hop --- gone in both.

\begin{table}[H]\centering\small
\begin{tabular}{clll}
\toprule
step & type & \armA{} best (0.843) & \armB{} best (0.858) \\
\midrule
1 & \textcolor{accent}{tool} & \texttt{retrieve\_deep(\$outer)} & \texttt{retrieve\_deep(\$outer)} \\
2 & \textcolor{bad!70!black}{llm} & 2nd-hop query \hfill $[1]$ & 2nd-hop query \hfill $[1]$ \\
3 & \textcolor{accent}{tool} & \texttt{retrieve\_deep(\$hist[-1])} \hfill $[2]$ & \texttt{retrieve\_deep(\$hist[1])} \hfill $[2]$ \\
4 & \textcolor{bad!70!black}{llm} & 3rd-hop query \hfill $[1,3]$ & 3rd-hop query \hfill $[1,3]$ \\
5 & \textcolor{accent}{tool} & \texttt{retrieve\_deep(\$hist[-1])} \hfill $[4]$ & \texttt{retrieve\_deep(\$hist[3])} \hfill $[4]$ \\
6 & \textcolor{bad!70!black}{llm} & 4th-hop query \hfill $[1,5]$ & 4th-hop query \hfill $[1,3,5]$ \\
7 & \textcolor{accent}{tool} & \texttt{retrieve\_deep(\$hist[-1])} \hfill $[6]$ & \texttt{retrieve\_deep(\$hist[5])} \hfill $[6]$ \\
\bottomrule
\end{tabular}
\caption{Convergent evolution, independently in both arms: identical step types, identical
topology, near-identical dependency sets. The differences are cosmetic --- \armA{} names a tool
step's query source relatively (\texttt{\$history[-1]}), \armB{} absolutely
(\texttt{\$history[1]/[3]/[5]}), functionally equivalent on a linear chain --- plus one dropped
skip-edge in \armA's step 6 ($[1,5]$ vs $[1,3,5]$). Both operators find ``deepen every hop, drop
the summariser.'' (\texttt{\$outer} = the claim under verification; \texttt{\$history[k]} = the
output of step $k$, \texttt{[-1]} = the previous step.)}
\end{table}

\subsection{How \armB{} evolved it: the base\,$\to$\,best lineage}
\label{sec:lineage}
The best \armB{} child sits at the end of a clean four-mutation lineage from the seed. Every
child is a structured slot-diff of its parent, so the whole path is legible as a sequence of
typed edits --- a trail the free-rewrite arm lacks, since each of its children is an independent
whole-document re-emission with no operator-native diff (its deltas would have to be
reconstructed post hoc).

Read Fig.~\ref{fig:lineage} left to right and the graph is the story. \textbf{gen2} fans
skip-edges into the two query steps (\#4,\,\#6) so each hop's query sees the whole accumulated
history --- the decisive edit (train $0.753 \to 0.847$), a coordinated seven-step rewrite in
which every tool step also becomes \texttt{retrieve\_deep} and the summarise steps flip to
query-generators. \textbf{gen3} strips those fan-ins back to a clean linear chain while
rewriting the reasoning prompts ($+0.009$). \textbf{gen4} re-adds a selective subset of
skip-connections for the final $+0.002$. base\,$\equiv$\,gen1 carries no change at all: the
first mutation was a representable no-op. Because each child is a diff against a known parent,
the model touches only the wiring it means to --- never re-emitting (and re-risking) the
untouched steps.

\begin{figure}[p]\centering
\includegraphics[width=\textwidth,height=0.92\textheight,keepaspectratio]{lineage_graph.png}
\caption{\armB's winning lineage as a \textbf{dependency graph}, one column per program
(base\,$\equiv$\,gen1\,$\to$\,gen4). Each column is the seven-step pipeline drawn top-to-bottom;
every node is a step, colored by type (\textcolor{accent}{blue}~=~\texttt{retrieve}/%
\texttt{retrieve\_deep} tool, \textcolor{bad!70!black}{orange}~=~LLM) and labelled with its role.
\textbf{Arrows are the step \texttt{dependencies}} --- the wiring that defines the DAG. An edge is
grey when it is unchanged from the parent, \textcolor{good}{solid green} when the mutation
\emph{added} it, and \textcolor{bad}{dashed red} when it \emph{removed} it; a node is outlined in
\textcolor{bad}{red} when any of its fields changed versus the previous generation. The italic line
under each new generation is the CARL structured-diff operator's own stated change --- its words,
not a reconstruction.}
\label{fig:lineage}
\end{figure}

\begin{table}[H]\centering\footnotesize
\begin{tabular}{llccc}
\toprule
program & structural change vs.\ parent & fitness (train soft) & test soft & test disc.\ (\% $\pm$ SE) \\
\midrule
base (seed)     & ---                              & ---   & 0.752 & 44.0 $\pm$ 2.9 \\
gen 1           & no-op (structurally identical)   & 0.753 & 0.763 & 47.7 $\pm$ 2.9 \\
gen 2           & rewire every hop; \texttt{+}skip-edges & 0.847 & 0.848 & 64.3 $\pm$ 2.8 \\
gen 3           & prune fan-ins to linear chain    & 0.856 & 0.859 & 66.7 $\pm$ 2.7 \\
gen 4 (best)    & re-add selective skip-edges      & \textbf{0.858} & \textbf{0.862} & \textbf{\textcolor{good}{68.0 $\pm$ 2.7}} \\
\midrule
\armA{} best (free rewrite) & --- (no inspectable diff) & 0.843 & 0.853 & 67.3 $\pm$ 2.7 \\
\bottomrule
\end{tabular}
\caption{Held-out test payoff of \emph{every} program on the lineage ($n=300$ test claims; fitness
is the train soft-coverage objective, 300 train claims). The $\pm$ on discrete coverage is the
binomial claim-sampling SE, $\sqrt{p(1-p)/300}$; it captures test-set sampling only --- executor
re-draw noise comes on top. Per-claim soft scores were not persisted, so the soft columns carry
point values only. Held-out discrete coverage climbs
monotonically down the four-mutation path ($44.0\to47.7\to64.3\to66.7\to\textbf{68.0}$) in this
single draw, soft coverage tracking in lockstep; only the gen1$\to$gen2 jump ($+16.6$ pts)
clearly exceeds the noise floor. base and gen1 are the \emph{same} structure yet score $44.0$ vs $47.7$: that
$\sim$3.7-pt gap is the non-deterministic \texttt{Qwen3-8B} executor's noise floor, the band
against which per-generation deltas --- and the final $68.0$ vs $67.3$ winner comparison ---
should be read. On held-out data the two arms' winners are a statistical tie.}
\label{tab:lineagetest}
\end{table}

\noindent Two structural observations. First, the operator can emit a \emph{no-op} diff
(\emph{base\,$\to$\,gen1}): keeping every slot unchanged is a representable, valid move --- wasted
compute, but never a malformed child. Second, the load-bearing mutation (\emph{gen1\,$\to$\,gen2})
is a \emph{coordinated multi-field, multi-step rewrite} --- tool swaps, step-type flips, and
dependency rewiring together --- yet it is expressed as one structured diff that the compiled
schema guarantees is well-formed. The very edit that is most likely to corrupt a free-form
re-emission (touch every step at once) is exactly where the diff's constraints pay off.

\partrule{Part II --- top-tier mutator: \texttt{Qwen3-235B-A22B-Thinking}}{full 250-mutant paired run --- does the diff's edge survive a frontier-scale mutator?}

\section{Results (Qwen mutator)}
\label{sec:resultsq}

\subsection{Validity funnel --- 100\% vs.\ a serving artifact}
\begin{table}[H]\centering
\begin{tabular}{lcc}
\toprule
full 250-mutant run & \armA{} (rewrite) & \armB{} (tool diff) \\
\midrule
mutation LLM calls (incl.\ retries) & 265 & 268 \\
pre-persist generation failures & 9 (no code extracted) & 17 (call errors; \textbf{0 schema}) \\
children persisted & 250 & 250 \\
\quad still executing at stop (excluded) & 4 & 5 \\
\midrule
\textbf{completed evaluations} & \textbf{246} & \textbf{245} \\
\quad valid & 211 & \textbf{245} \\
\quad invalid & 35 & \textbf{\textcolor{good}{0}} \\
\qquad of which malformed wire-JSON & \textcolor{bad}{34 (13.8\%)} & \textbf{\textcolor{good}{0}} \\
\qquad other (well-formed, failed validation) & 1 & 0 \\
\textbf{valid yield among completed} & 85.8\% & \textbf{\textcolor{good}{100.0\%}} \\
\bottomrule
\end{tabular}
\caption{On completed evaluations the diff arm is clean: 245 of 245 valid, zero malformed,
and --- notably --- \textbf{zero schema rejections in 268 calls} (the gemini pair had 2), so the
compiled grammar never even had to bounce a payload. Operational waste remains upstream: 17
pre-persist failures, transient LLM-call errors (proxy timeouts on the thinking model), each
costing one retried call, nothing persisted. \armA's 34 malformed children look like a
\emph{serving artifact}, not model confusion: every one is \emph{head-truncated} --- the
document begins mid-stream (typically at \texttt{system\_prompt}, the opening \texttt{\{"}
clipped) --- consistent with a structured-output defect of the local serving stack under long
thinking-model emissions. \armA's true content-level failure rate is \textbf{1/246 = 0.4\%}: one
well-formed document that failed chain validation. (Funnel rows again do not close exactly:
$265-9=256$ vs 250 persisted in \armA{} --- six calls unclassifiable from the logs --- and
$268-17=251$ vs 250 in \armB.)}
\end{table}
\begin{figure}[H]\centering
\includegraphics[width=\textwidth]{ab_overview_qwen.png}
\caption{Left: per-child fitness and best-so-far. \armA's invalid crosses dot the floor line
(the head-truncation artifact); \armB's floor is empty. \armB{} jumps early (its winner exists
by rank 60) and plateaus; \armA{} grinds upward and posts its best near the very end (rank 236).
Right: the validity funnel --- \armB{} keeps every completed child.}
\end{figure}

\noindent\textbf{The same lesson from the other end of the document.} Part~I's residual
free-rewrite failures were tail-truncated (OpenRouter/gemini); Part~II's are head-truncated
(local vLLM-style stack under a thinking model). Different providers, different ends, same
class: \emph{a whole-document rewrite is one long fragile emission, and each of the two stacks
we ran on found its own way to cut it}. The diff's slot-local payloads plus provider-side
grammar were untouched on both --- zero persisted malformed children in either diff run; and
had an outer diff response been truncated, the grammar-checked parse would have failed
pre-persist rather than persisting a corrupt genome.

\subsection{Fitness --- a tie at the top, wider dispersion below}
\begin{table}[H]\centering
\begin{tabular}{lcc}
\toprule
completed children & \armA{} (rewrite) & \armB{} (tool diff) \\
\midrule
best child & \textbf{0.824} & 0.819 \\
child mean $\pm$ SD & $0.757 \pm 0.070$ & $0.720 \pm 0.129$ \\
child median & 0.768 & 0.758 \\
seed (executor re-score) & 0.734 & 0.743 \\
\bottomrule
\end{tabular}
\caption{Best fitness is a tie: the 0.005 gap is a quarter of the 0.021 seed re-score band ---
and it is \armA{} that noses ahead this time, despite \emph{fewer} valid draws (211 vs 245), so
best-of-$N$ bias cannot explain it. \armA's winner arrives at chronological rank 236/250; \armB's
at rank 60, never beaten for the remaining 190 children. \armB's lower mean is a
\emph{dispersion} story, not short chains alone: restricted to 7-step children the gap persists
(\armA{} $0.766 \pm 0.044$, $n{=}171$ vs \armB{} $0.735 \pm 0.131$, $n{=}175$) --- \armB{}
takes bolder edits with a heavy low tail at matched topology, where its median (0.768) stays
close to \armA's (0.773). The 70 \armB{} children below 7 steps (mean 0.682;
\S\ref{sec:structq}) add a further drag on the mean.}
\end{table}

\subsection{Held-out test coverage}
\label{sec:benchcmpq}
Same protocol as Part~I: $n=300$ held-out HoVer claims, discrete coverage $=1$ iff every gold
title is retrieved, plus soft coverage. Programs evaluated: the seed (fresh re-score), both
winners, and \armB's decisive intermediate generation.

Part~II generalises less dramatically than Part~I. Both winners clear the seed ($48.3 \pm 2.9$)
--- \armA{} $59.0 \pm 2.8$, \armB{} $56.7 \pm 2.9$ (binomial claim-sampling SE, $n=300$) --- but
the margins over the prior sweep's without-feedback best
(56.0, Table~\ref{tab:benchcmp}) are 0.7--3.0 points, inside the $\sim$3.7-pt single-draw noise
floor; there is no Part~I-style 11-point clearance here. The held-out ranking also scrambles the
train ranking: \armB's \emph{intermediate} gen-2 posts the best held-out score of the whole pair
(61.0) while its train-fitter child gen-3 drops back to 56.7. Two readings, both honest: between
the arms, held-out rank is a coin flip at this noise level; and train gains past gen-2 did not
transfer --- consistent with the final prompt-only tweak fitting the train draw, or simply with
noise. Soft coverage tells the same story in miniature (seed 0.767; all three evolved programs
0.816--0.833).

\begin{table}[H]\centering\small
\begin{tabular}{lccc}
\toprule
program & train soft & test soft & test disc.\ (\% $\pm$ SE) \\
\midrule
base (seed)                     & ---            & \qTseedS & \qTseedD{} $\pm$ 2.9 \\
\armB{} gen 2 (rewire)          & 0.804          & \qTgS    & \qTgD{} $\pm$ 2.8 \\
\armB{} gen 3 (best)            & 0.819          & \qTbS    & \qTbD{} $\pm$ 2.9 \\
\armA{} best (free rewrite)     & \textbf{0.824} & \qTaS    & \qTaD{} $\pm$ 2.8 \\
\bottomrule
\end{tabular}
\caption{Held-out generalisation of the Part~II winners and \armB's lineage ($n=300$). $\pm$ is
the binomial claim-sampling SE, as in Table~\ref{tab:lineagetest}; soft columns are point values
(per-claim scores not persisted). The evolved programs' lifts over the seed clear 2--3~SE of the
difference, but every gap \emph{between} evolved programs is $\lesssim$1~SE --- the held-out ranking
is unresolved at this sample size. These rows also fill the bottom block of
Table~\ref{tab:benchcmp} --- the closest available comparison against the prior
feedback-condition sweep under the same mutator.}
\label{tab:lineagetestq}
\end{table}

\subsection{Token cost --- no reasoning split on the proxy}
\label{sec:tokensq}
The proxy serving \texttt{Qwen3-235B-Thinking} reports total completion tokens only --- the
thinking channel is not itemised on either arm (the split below is structurally unavailable,
unlike Part~I where it is now fully measured). Totals are still cross-arm comparable: output
\emph{includes} thinking on both arms.
\begin{table}[H]\centering
\begin{tabular}{lccc}
\toprule
mutation side, full run (Qwen-235B) & \armA{} & \armB{} & \armB{} vs \armA{} \\
\midrule
tokens in & 1{,}512{,}046 & 1{,}424{,}748 & $-5.8\%$ \\
tokens out (total, incl.\ thinking) & 2{,}278{,}810 & 2{,}629{,}263 & $+15.4\%$ \\
LLM calls (incl.\ retries) & 265 & 268 & $+1\%$ \\
\bottomrule
\end{tabular}
\caption{The thinking model generates more than it reads on both arms. \armB{} spends 15\% more
output --- on a thinking model the diff cannot save total output, because reasoning (the bulk)
scales with the \emph{decision}, not the payload; Part~I's measured split (79--87\% reasoning)
says the payload economy is real but small against the thinking channel. \armB{} reads 6\% less:
the schema-bearing prompt is cheaper than \armA's full-document-plus-instructions prompt at this
tokenizer.}
\end{table}
\begin{figure}[H]\centering
\includegraphics[width=\textwidth]{ab_tokens_qwen.png}
\caption{Left: output tokens per mutation call (10-call rolling mean) --- \armB{} runs
consistently hotter, its extra thinking driven by the slot-choice decision. Right: cumulative
input/output over the run.}
\end{figure}

\subsection{Cost/quality frontier --- the winner for a quarter of the budget}
\label{sec:paretoq}
Same attribution as Part~I (\texttt{created\_at} ordering), priced at the cheapest OpenRouter
rate for \texttt{Qwen3-235B-A22B-Thinking} (\$0.1495/M in, \$1.495/M out) for comparability ---
the actual serving is self-hosted.

\begin{table}[H]\centering
\begin{tabular}{lccc}
\toprule
mutation-model \$ (OpenRouter cheapest) & \armA{} & \armB{} & \armB{} vs \armA{} \\
\midrule
total run spend & \$3.63 & \$4.14 & $+14\%$ \\
\midrule
\textbf{cost to the winning program} & \$3.42 & \textbf{\textcolor{good}{\$1.02}} & \textbf{$-70\%$} \\
\quad winner reached at (chrono.\ rank) & 236/250 (94\%) & \textbf{\textcolor{good}{60/250 (24\%)}} & far earlier \\
\quad winner fitness (train soft) & 0.824 & 0.819 & tie (noise band) \\
\midrule
cost per \emph{valid} program & \$0.0172 & \$0.0169 & parity \\
budget wasted on invalid children & \textcolor{bad}{\$0.52 (14.2\%)} & \textbf{\textcolor{good}{\$0.00}} & --- \\
\bottomrule
\end{tabular}
\caption{The retrospective timing gap widens at the top end: \armB's never-beaten winner exists
after 24\% of the run for \textbf{\$1.02}; \armA's does not exist until 94\% --- stopping
\armA{} early at \armB's \$1.02 mark would have forfeited nearly all its eventual gain.
Cost-per-valid is parity (\armA's 14\% waste on head-truncated children is offset by \armB's
heavier thinking). At the illustrative OpenRouter prices used throughout, the whole pair cost
\$7.78 against Part~I's \$39.59 --- about $5\times$ cheaper for comparable peaks; actual
economics depend on serving costs.}
\end{table}
\begin{figure}[H]\centering
\includegraphics[width=\textwidth]{ab_pareto_qwen.png}
\caption{\textbf{Left:} best-so-far train fitness vs.\ cumulative mutation dollars --- \armB{}
jumps to its plateau by \$1.02; \armA{} climbs in steps and crosses only at the very end of its
budget. \textbf{Right:} cost-to-program vs.\ held-out test discrete coverage for the evaluated
programs of Table~\ref{tab:lineagetestq}.}
\end{figure}
\begin{figure}[H]\centering
\includegraphics[width=\textwidth]{ab_waste_qwen.png}
\caption{\textbf{Left:} budget split into productive and wasted spend --- \armA{} burns \$0.52
(14.2\%) on head-truncated children; \armB{} wastes \$0.00. \textbf{Right:} cost per valid
program --- parity at $\approx$\$0.017, the waste offset by thinking-token overhead.}
\end{figure}

\subsection{Structural exploration --- the diff explores \emph{wider} under a strong mutator}
\label{sec:structq}
Part~I's pattern inverts. \armA's valid children stay near the seed topology:
$\{7\!:\!171,\ 6\!:\!27,\ 5\!:\!12,\ 4\!:\!1\}$ (81\% at 7 steps, nothing below 4). \armB's
spread across the whole admissible range:
$\{7\!:\!175,\ 6\!:\!36,\ 5\!:\!16,\ 4\!:\!11,\ 3\!:\!3,\ 2\!:\!4\}$ --- 34 children at 2--5
steps, including deliberate 2-step probes. With a mutator strong enough to reason about
topology, the typed slot grammar functioned here as an \emph{exploration} instrument: pruning to
a 4-slot chain is one legal diff, whereas the free rewriter, re-emitting everything, gravitated
to reproducing the shape it was shown. Those short probes account for part --- not all; see the
dispersion note above --- of \armB's lower child mean; the return is topology coverage, at no
cost to the peak.
\begin{figure}[H]\centering
\includegraphics[width=0.72\textwidth]{ab_nsteps_qwen.png}
\caption{Chain-length distribution of valid children under the Qwen mutator. \armB{} (blue)
covers 2--7 steps; \armA{} (red) never goes below 4 and concentrates at the seed's 7.}
\end{figure}

\section{The winning chains (Qwen) --- divergent architectures}
\label{sec:chainsq}
Unlike Part~I, the arms do \emph{not} converge. Under the frontier mutator each operator found a
genuinely different retrieval strategy --- beating Part~II's seed by 0.09 (\armA) and 0.08
(\armB) train soft:

\begin{table}[H]\centering\footnotesize
\begin{tabular}{cllll}
\toprule
 & \multicolumn{2}{c}{\armA{} best (0.824, gen 9)} & \multicolumn{2}{c}{\armB{} best (0.819, gen 3)} \\
\cmidrule(lr){2-3}\cmidrule(lr){4-5}
step & type & step content & type & step content \\
\midrule
1 & \textcolor{accent}{tool} & \texttt{retrieve\_deep(\$outer)} & \textcolor{accent}{tool} & \texttt{retrieve(\$outer)} \\
2 & \textcolor{bad!70!black}{llm} & 2nd-hop query, branch 1 \hfill $[1]$ & \textcolor{bad!70!black}{llm} & 2nd-hop query \hfill $[\,]$ \\
3 & \textcolor{accent}{tool} & \texttt{retrieve\_deep(\$hist[-1])} \hfill $[2]$ & \textcolor{accent}{tool} & \texttt{retrieve(\$hist[1])} \hfill $[2]$ \\
4 & \textcolor{bad!70!black}{llm} & 2nd-hop query, branch 2 \hfill $[1,2,3]$ & \textcolor{bad!70!black}{llm} & ``summarise'' $\to$ query \hfill $[\,]$ \\
5 & \textcolor{accent}{tool} & \texttt{retrieve\_deep(\$hist[-1])} \hfill $[4]$ & \textcolor{bad!70!black}{llm} & 3rd-hop query \hfill $[\,]$ \\
6 & \textcolor{bad!70!black}{llm} & combine $\to$ 3rd-hop query \hfill $[1,3,5]$ & \textcolor{accent}{tool} & \texttt{retrieve\_deep(\$hist[4])} \hfill $[5]$ \\
7 & \textcolor{accent}{tool} & \texttt{retrieve\_deep(\$hist[-1])} \hfill $[6]$ & \textcolor{accent}{tool} & \texttt{retrieve\_deep(\$hist[3])} \hfill $[4]$ \\
\bottomrule
\end{tabular}
\caption{Two different answers to the same task. \armA's winner is \textbf{evidence-conditioned
branch-and-merge}: two parallel second-hop query branches, each reading the first-hop passages,
merged by a combiner that issues the third hop --- all \texttt{retrieve\_deep}. \armB's winner is
\textbf{claim decomposition}: three query-generating LLM steps whose dependency sets are
\emph{empty} (they read only the claim), feeding four retrievals, the last two
\texttt{retrieve\_deep} in \emph{parallel} terminal position. Since fitness is the union of gold
titles over all tool outputs, parallel claim-conditioned probes are a legitimate strategy ---
query \emph{diversity} substitutes for evidence chaining. Note the vestigial title: \armB's step
4 still says ``summarise second-hop evidence'' but its inputs were cut by the gen-2 rewire; it
now functions as another claim-conditioned query generator whose output step 7 retrieves.}
\end{table}

\subsection{How \armB{} evolved it: two mutations, the second prompt-only}
\label{sec:lineageq}
\armB's winner is only \textbf{two diffs} from the seed, and the operator's own logged
justifications narrate the search:
\begin{itemize}
  \item \textbf{seed $\to$ gen 2} ($0.743 \to 0.804$, $\Delta +0.061$): a coordinated rewire.
  The operator's stated change: \emph{``Removed first summarization step and added fourth
  retrieval hop using two-hop summary: the summarization step discarded critical evidence
  fragments needed for multi-hop reasoning, and the fourth hop recovers missing gold articles by
  re-querying''} --- acting, per its justification, on a ranked program insight against the lossy
  first-hop summariser. One diff: drop a step, add a fourth tool step, retype two tool steps to
  \texttt{retrieve\_deep}, rewire the query sources.
  \item \textbf{gen 2 $\to$ gen 3 (best)} ($0.804 \to 0.819$): \textbf{a
  single-field, prompt-only edit}. Mutation archetype (the operator's self-selected edit
  category): \emph{Precision Optimization}; justification:
  \emph{``the program is at the frontier but the window trend is flat \dots without altering the
  proven topology, as previous structural changes regressed.''} The whole diff is one
  \texttt{keep} slot gaining a \texttt{reasoning\_questions} field: \emph{``What specific fact is
  missing to verify the claim? How can this fact be expressed as a search query?''} Every other
  slot: \texttt{keep}, untouched.
\end{itemize}
That last mutation is the legibility argument in miniature: the winning edit of a 250-child run
under a frontier mutator is \emph{readable verbatim} --- one prompt field, two short questions,
stated intent --- because the operator emits typed diffs. The free-rewrite arm's gen-9 winner
offers no such trail: nine whole-document re-emissions whose actual deltas would have to be
reverse-engineered.

\begin{figure}[p]\centering
\includegraphics[width=\textwidth,height=0.92\textheight,keepaspectratio]{lineage_graph_qwen.png}
\caption{\armB's Part~II winning lineage (base $\to$ gen2 $\to$ gen3), same visual grammar as
Fig.~\ref{fig:lineage}: columns are programs, nodes are steps (\textcolor{accent}{blue} tool /
\textcolor{bad!70!black}{orange} LLM), \textcolor{good}{green} edges added, \textcolor{bad}{red
dashed} removed, red outline = fields changed. The gen-2 column shows the rewire: the first
summariser gone, a fourth retrieval added terminally, the LLM steps' evidence edges cut ---
turning the chain from sequential evidence-chaining into parallel claim-decomposition. The gen-3
column is visually almost identical to gen-2 --- a single red-outlined node (step 2) with no
wiring change: the prompt-only mutation described above.}
\label{fig:lineageq}
\end{figure}

\section{What the first arm-A runs measured: structured-output harness hygiene}
\label{sec:artifact}
This report's arm-A runs are re-runs. The original pair produced the intended headline ---
\emph{``even a strong mutator emits malformed JSON in 38\% of free rewrites''} --- and it did not survive an
audit. The episode is worth a section because the failure mode is general to any harness that
pairs a prompt with a structured-output schema.

\subsection{The bug}
The shared mutation agent's structured-output model carried a field description written for
Python-genome problems: \texttt{MutationStructuredOutput.code} was described as \emph{``complete
mutated Python source code \dots NEVER put JSON''} --- while the HoVer problem prompt instructed
\emph{``\texttt{code} is one COMPLETE wire-format JSON document --- not Python.''} Two
instructions at war inside one call. Gemini followed the schema description (emitting Python) in
$\sim$36\% of calls; Qwen, notably, never did --- its pre-fix penalties were already the
head-truncation serving artifact of \S\ref{sec:resultsq}, misfiled under the same error label.

\subsection{Decomposition of the ``malformed'' children}
Content inspection of every penalty child (a child whose evaluation drew the $-1000$ catch-all
penalty fitness) in the original runs:
\begin{itemize}
  \item \textbf{gemini (armA, 93 penalties):} 87 were \emph{syntactically valid Python}
  (\texttt{def entrypoint():\,\dots}) --- the model doing exactly what the schema description
  asked --- plus 2 truncations, 1 bad escape, 1 control character, 2 well-formed-but-invalid.
  True document-syntax breakage: \textbf{4/242 $\approx$ 1.7\%}, not 38\%.
  \item \textbf{Qwen (armA, 29 penalties, partial run):} 28 head-truncated fragments and 1
  well-formed-but-invalid --- i.e.\ Qwen's failures were \emph{never} the schema contradiction;
  they were already the serving artifact that persists in the fixed re-run
  (\S\ref{sec:resultsq}).
\end{itemize}
The fix (a format-agnostic \texttt{code} description) plus full re-runs moved gemini's malformed
rate $38\% \to 3.3\%$ --- and left Qwen's essentially unchanged ($\sim$12\% $\to$ 13.8\%),
confirming the two models were failing for entirely different reasons that a single
``malformed'' error kind had collapsed into one number.

\begin{eli5}
The mutator was handed a form whose cover letter said ``answer in JSON'' while the fine print on
the answer box said ``answer in Python, never JSON.'' Some models trust the cover letter, some
the fine print --- and everything written in Python got stamped MALFORMED by a grader expecting
JSON. The fix was to make the fine print agree with the cover letter. Arm~B never had this
problem: its answer box \emph{physically} only accepts one format --- there is no fine print to
disagree with.
\end{eli5}

\subsection{Lessons}
\begin{enumerate}
  \item \textbf{Schema descriptions are prompts.} A structured-output field description competes
  with the instruction text for the model's obedience; audit both sides of every call.
  \item \textbf{Taxonomies need content inspection.} The error-kind label said
  \texttt{json\_parse\_error} for all 93 children; the contents said ``87 valid Python.'' A
  penalty bucket is a claim about the \emph{harness} until the payloads are read. (Likewise,
  \armB's 3 gemini ``invalid'' children turned out to be evaluator crashes with perfectly
  parseable genomes --- a $-1000$ catch-all conflates genome defects with infra failures.)
  \item \textbf{Grammar-side constraints are immune to this class.} \armB's numbers are
  identical before and after the fix, because under a compiled \texttt{json\_schema} grammar
  the contradiction never gets a vote: the fine print cannot override what is sampleable. The
  diff's zero-malformed record thus covers not only model error but \emph{harness} error --- the
  category this section exists to document.
\end{enumerate}

\section{Cross-model synthesis --- what shrinks, what holds}
\label{sec:synthesis}
The planned ``malformed-JSON ladder'' (weak llama 40\% $\to$ gemini 38\% $\to$ Qwen $\sim$12\%
$\to$ schema 0\%) died in the audit: its middle rungs were measuring a harness bug and a serving
artifact, not model JSON competence. What replaces it is sharper:

\begin{table}[H]\centering\footnotesize
\begin{tabular}{llll}
\toprule
arm-A failure class & Part~I (gemini) & Part~II (Qwen-235B) & \armB{} (both) \\
\midrule
harness artifact (schema/prompt war) & 38\% $\to$ \textbf{0} after fix & never occurred & unrepresentable \\
truncated emission (attributed to serving stack) & 3.3\% (tail-cut) & 13.8\% (head-cut) & none persisted \\
true content defect (well-formed, invalid) & $\sim$0 & 0.4\% & 0 \\
\bottomrule
\end{tabular}
\caption{Where free-rewrite children actually die, by cause. Model IQ was never the story: with
the harness honest, \emph{neither} strong mutator meaningfully miswrites JSON. What both stacks
do is \emph{truncate long single-string emissions} --- from opposite ends --- and only the
whole-document rewrite ships a payload long enough to be at risk. The diff is robust not just to
model error but to the prompt bug and to both serving stacks' quirks: \textbf{zero persisted
malformed children in both diff runs}.}
\end{table}

\noindent\textbf{What shrinks with mutator strength.} The fitness margin: Part~I leans \armB{}
(+0.015 best), Part~II leans \armA{} (+0.005) --- both inside the 0.021 same-genome re-score
band. So does the waste margin as a share of spend that \emph{matters}: cost-per-valid-program
is parity in both pairs (\$0.084 vs \$0.083; \$0.0172 vs \$0.0169). A strong mutator makes the
free rewrite \emph{survivable}: you lose 3--14\% of children to truncation and match the diff's
peak.

\medskip\noindent\textbf{What holds at every strength.}
\begin{itemize}
  \item \textbf{Validity of persisted genomes is categorical.} 0 malformed for the diff in both
  its runs; in Part~II, 0 invalid of any kind and 0 schema rejections in 268 calls.
  (Operational waste --- transient call errors, one no-op diff --- remains.)
  \item \textbf{The winner arrives early (retrospectively).} Chronological rank 80/254 (31\%)
  and 60/250 (24\%) for the diff, vs 131/252 (52\%) and 236/250 (94\%) for the rewrite;
  cost-to-winner \textbf{$-36\%$} and \textbf{$-70\%$}. In these runs, constrained local edits
  climbed fast from a working parent and plateaued; whole-document rewriting drifted toward its
  peak, in Part~II finding it only in the final 6\% of the budget.
  \item \textbf{Legibility.} Part~I's winning lineage reads as four typed edits with stated
  intent; Part~II's winner is two diffs, the second a \emph{single prompt field}. The rewrite
  arms' winners (gen 4 and gen 9) have no operator-native path; their deltas would have to be
  reconstructed post hoc from whole-document emissions.
  \item \textbf{Structural range.} Under the strong mutator the diff explored 2--7-step
  topologies while the rewrite stayed at 4--7 hugging the seed's shape --- the grammar
  functioned as an exploration instrument here, at a measurable mean-fitness dispersion cost and
  no cost to the peak.
\end{itemize}

\medskip\noindent\textbf{Two asides the pair design surfaced.} First, architecture: the two
Part~I winners converged on the same deepen-and-accumulate chain, but the two Part~II winners
solved retrieval \emph{differently} (evidence-conditioned branching vs.\ claim decomposition) ---
in this pair, operator choice co-varied with \emph{which} solution family the search found, not
just how safely it got there (one replicate; suggestive, not established). Second, economics: the self-hosted frontier mutator delivered
comparable peaks at OpenRouter-priced \$7.78 for the whole pair vs.\ \$39.59 for gemini ---
at these illustrative prices, reasoning-token rates rather than capability drove the bill;
actual serving economics differ.

\section{Engineering deliverable: per-candidate executor token accounting}
\label{sec:tokentrack}
These runs carry \emph{no} per-candidate executor tokens; the only executor figure available is
a shared-proxy aggregate ($\approx$132M tokens for the Part~I pair, proxy Prometheus counter, no
per-arm split). The cause: the chain executor (\texttt{Qwen3-8B}) already captured per-call
\texttt{prompt\_tokens}/\texttt{completion\_tokens} in its \texttt{CallLog}, but the chain
runners call \texttt{client.copy()} for every concurrent LLM call and the base \texttt{copy()}
hands out a \emph{fresh} log list for thread-safety --- so usage landed in throwaway copies, the
parent's logs stayed empty, and \texttt{validate()} discarded it.

That gap is now closed. \texttt{problems/chains/usage.py} adds a \texttt{LogAggregatingLLMClient}
whose \texttt{copy()} shares one log list across the fan-out, plus \texttt{usage\_totals()}
(prompt/completion/total tokens, \texttt{n\_llm\_calls}, \texttt{llm\_cost}) and a
\texttt{ZERO\_USAGE} constant. It is wired into \texttt{full7/validate.py} (success and
early-reject paths) and refactors the summarizer validator onto the same helper; four unit tests
cover the copy-sharing invariant. The same maintenance pass fixed the mutation-side
instrumentation: \texttt{tokens\_reasoning} is now emitted on \emph{all} \texttt{LLM\_CALL}
paths (it was silently dropped by the stock mutation agent --- the gap that blinded the pre-fix
token analysis), which is why Part~I carries a full reasoning split. \textbf{Future} HoVer runs
will additionally carry per-candidate executor tokens in program metrics, visible via
\texttt{gigaevo top} and the A/B analyzer.

\section{Verdict}
\label{sec:verdict}
\textbf{Q1 --- does the diff generalise to tool-bearing genomes?} Yes. In both of its runs the
tool-aware grammar persisted zero malformed children, edited tool steps as first-class slots
(retyping \texttt{retrieve} $\to$ \texttt{retrieve\_deep}, moving query sources, rewiring
dependencies), and in Part~II never triggered a single schema rejection in 268 calls. Residual
operational waste (transient call errors, one no-op diff) is disclosed above.

\medskip\noindent\textbf{Q2/Q3 --- does the guarantee matter under strong mutators?} Not for the
reason we expected. The naive version --- ``strong models still can't emit JSON'' --- is dead:
that number was our own harness bug, and this report documents it rather than buries it
(\S\ref{sec:artifact}). The considered version survives: a whole-document rewrite is one long
emission, and \emph{both} serving stacks truncated a slice of them (3.3\% tail-cut, 13.8\%
head-cut) --- a tax that did not depend on model IQ and that the diff's short, grammar-bound
payloads avoided on both stacks. In these runs the guarantee covered the whole stack: model,
prompt, \emph{and} provider.

\medskip\noindent\textbf{The surviving economic case is timing, not waste.} At parity
cost-per-valid-program, the diff's never-beaten winner existed after 31\% (Part~I) and 24\%
(Part~II) of the budget --- retrospectively $-36\%$ and $-70\%$ cost-to-winner --- while free
rewriting needed 52\% and 94\%. Banking that saving online needs a stopping rule; for a search
loop with any early-stopping or successive-halving discipline, it is the operative number. Peaks are statistical ties in both pairs.

\medskip\noindent\textbf{Recommendation.} Default to the structured diff for wire-JSON genomes
at both mutator strengths tested, under the conditions studied here (single-parent, memory-off,
HoVer retrieval chains). Under weak mutators (prior report) it is the difference between a
search and a syntax lottery; under strong ones it buys earlier winners, zero persisted malformed
output, wider topology exploration, and lineages you can read --- including, in Part~II, a
winning mutation that is one auditable prompt field. Keep free rewrite as an occasional
exploration control: in Part~II it found a genuinely different (branching) architecture and the
nominal best-fitness, so it still earns a lane as a diversity source where its truncation tax is
affordable.

\medskip\noindent\textbf{Caveats.} One replicate per arm per pair; no same-config variance floor
beyond the four seed re-scores (band 0.021); \armB{} runs predate the fixed \armA{} runs
(identical configs, non-concurrent execution; shared proxy and cluster load uncontrolled); fitness and held-out scores are single stochastic
draws (executor $T{=}0.6$; base-vs-gen1 puts the held-out noise floor near 3.7 pts); the $\pm$
on held-out discrete values is binomial claim-sampling SE only ($n=300$), with that re-draw
noise on top, and per-claim soft scores were not persisted (no $\pm$ on soft columns);
\texttt{memory=none}, \texttt{num\_parents=1} (no crossover); 4--12 in-flight children per arm
excluded at stop; Part~II reasoning/content split unavailable (proxy does not itemise thinking);
dollar figures are OpenRouter-priced tokens, executor excluded; a few mutation calls per run
are unclassifiable in the logs (funnel rows need not close exactly).

\end{document}
